\chapter{股票：把一家公司切成一亿片}
\chaptermeta{11}

\begin{ledetext}
买一股，到底买到了什么？不是一张纸，也不是一个代码，是排在所有人后面的那一份剩饭。
\end{ledetext}

\section{八月那张单子，钱是排队发的}

陈晓婷 32 岁，在杭州拱墅区开了家叫"三分甜"的奶茶店。八月结账，收银系统吐出来的营业额是 12 万 6400。

这笔钱不是她的。至少现在还不是。

\begin{flowlist}
  \flowitem{01}{供应商先拿}{珍珠、茶叶、鲜奶、杯子和封口膜，这个月进货 4 万 1200。不结账，下个月不发货。}
  \flowitem{02}{然后是四个店员}{工资加社保 3 万 4700。晚发三天，暑假工第二天就不来了。}
  \flowitem{03}{房东}{1 万 9600，押一付三，晚交按天算滞纳金。}
  \flowitem{04}{水电、外卖平台抽成、杂七杂八}{1 万 800。这一项每年都在往上爬，她已经不太看了。}
  \flowitem{05}{银行}{2019 年装修借的 42 万经营贷，这个月还息 1347。不还，抵押的那套小两居就不是她的了。}
  \flowitem{06}{税务局}{增值税、附加、企业所得税。有小微优惠，但不是不交。}
  \flowitem{07}{最后轮到陈晓婷}{1 万 7800 出头。去年八月对面开了家新店，这个数字是负的。}
\end{flowlist}

这个顺序不是行业习惯，是法律。公司清算的时候，职工工资和社保、税款、有担保的债权、普通债权，一层一层排下去，股东永远在最后一位。

法律给这个位置起了个名字：\term{剩余索取权}{residual claim}\index{048@剩余索取权}。买一股股票，买到的就是这个。

排最后有两个后果，是同一件事的两种长相。

坏的那面：前头任何一个人多拿一点，陈晓婷就少拿一点。营业额掉两成，房租工资利息一分不能少，那两成全从她这一份里扣。

好的那面：上不封顶。银行拿到 1347 就走了，供应商拿到货款就走了。哪个月这条街突然爆火，多出来的钱一分都不用分给他们。

\begin{pullquote}{}
股票风险高和股票长期回报高，不是两个事实，是同一个位置的两种读法。
\end{pullquote}

\section{王磊那 1800 块，三分甜一分钱没拿到}

三分甜后来开到了 14 家店，2025 年合并净利润 100 万，其中有三家还在亏钱。

这个体量当然不够上市。这本书里让它上市，是因为它的数字你能心算。

\begin{egbox}{举个栗子}
挂牌那天，三分甜发行 25 万股新股，每股 15 元，募到 375 万。这 375 万真金白银打进公司账户，用来在滨江开五家新店。这是\textbf{一级市场}，公司确实拿到了钱。

王磊在合肥一家汽车零部件厂做工艺工程师，月薪 1 万 3 出头。2021 年他拿了 3 万 8 的年终奖，抽出 1 万 2 开了证券账户。三分甜上市第二天开盘价 18 元，他觉得这生意他看得懂，买了 100 股，付出去 1800 块。

这 1800 块给了谁？给了把股票卖给他的那个人。三分甜的账户余额和昨天收盘时一模一样，一分没多。这是\textbf{二级市场}。

\end{egbox}

王磊以为自己在给公司投资。他其实是在跟另一个股民换手。

那二级市场是不是白长的？恰恰相反，只是它的作用全都隔了一层。

它给三分甜的未来标价。股价 18 元时再发 20 万股能拿到 360 万；跌到 6 元，同样 20 万股只能拿 120 万。融资成本明码写在屏幕上，谁都能看见。

它还是并购的钱。三分甜想吃掉隔壁的"茶小满"，可以不掏现金，直接拿自己的股票换。它也决定了店长手里那批行权价 10 元的期权值不值钱：股价 18 元，一股赚 8 块；股价 6 元，那就是一张彩色废纸，人留不住。

最要紧的是最后一条：它给一级市场开了个出口。

2019 年投三分甜第一笔钱的那位天使投资人，他的钱只有在能卖掉的那天才算落袋。没有二级市场，就不会有人做早期投资。风险投资、创业公司、期权激励，整条链子都吊在"将来能卖掉"这一个预期上。

把二级市场关掉，一级市场三个月内枯死。

\section{PE 是懒人版的贴现公式}

第 6 章说过，任何资产的价值，等于它未来所有现金流按某个折现率贴回今天。这个式子数学上无懈可击，实践中几乎没法用：要预测未来二十年的利润，还要挑一个折现率，手抖一下答案差三倍。

市场于是发明了一个偷懒的办法。

\begin{egbox}{举个栗子}
三分甜净利润 100 万，总股本 100 万股，每股收益 EPS = 1 元。王磊买入价 15 元，\term{市盈率}{P/E ratio}\index{052@市盈率} = 15 ÷ 1 = \textbf{15 倍}。

15 倍的字面意思：假设利润永远不变，公司要替他赚 15 年，才够把买入价还回来。

倒过来看更有味道。1 ÷ 15 = \textbf{6.7\%}，这是他的盈利收益率。同期 5 年期以上 LPR 是 3.5\%（2026 年 8 月），10 年期国债收益率更低。6.7\% 对 3.5\%，多出来那 3.2 个百分点，就是市场给"奶茶可能过气、店可能关门、陈晓婷可能不想干了"这一堆事开的价钱。

\end{egbox}

PE 把那个复杂的贴现公式压成了一个数字，代价是所有关于未来的假设全被塞进了分母里。它从来不是"高就贵、低就便宜"。

成长股 PE 高，买的是未来。一家利润每年涨 30\% 的公司，今年 PE 50 倍，三年后利润翻一倍多，股价不动的话 PE 只剩 22 倍。前提是那个未来真的来了。没来，50 倍就是 50 倍。

周期股 PE 低的时候，才是真要命的时候。

\begin{warnbox}{小心}
一家造集装箱的公司，行业景气顶点那年赚了 78 亿，股价对应 PE 只有 4.6 倍，看着比谁都便宜。

第二年运价崩了，利润掉到 8 亿多。股价一动没动，PE 立刻变成 45 倍。

\textbf{周期股 PE 最低的那一刻，通常正是利润最高的那一刻，也就是最危险的那一刻。}PE 的分母是当期利润，当期利润越不可持续，这个比值越骗人。这条规律翻过的车，能从上海排到杭州。

\end{warnbox}

\section{ROE 高不高，不如问它把钱放哪儿了}

ROE 是净利润除以股东权益，衡量股东投进去的每一块钱一年能赚回几毛。三分甜的净资产 500 万，净利 100 万，ROE 20\%。

它和增长之间有个朴素到近乎粗鲁的关系：

\textbf{可持续增长率 ≈ ROE ×（1 − 分红率）}

三分甜把 32\% 的利润分掉，留下 68\% 继续开店。它靠自己能长多快？20\% × 0.68 = 13.6\%。就这么多，除非去借钱或者增发。

所以"高 ROE"这三个字单独拿出来没有意义，得配一个问句：钱还有地方投吗？

三分甜每开一家新店都能赚回 20\%，那它就该少分红、多开店，这是一台复利机器在转。

可要是杭州的奶茶店已经饱和，新开一家只能赚 4\%，它还把利润死死留着扩张，等于把本来能赚 20\% 的钱拿去做 4\% 的生意。报表上写着"营收增长"，股东的钱在这个过程里被磨掉了。这种时候把钱还给股东，才是负责任的做法。

还钱有两种姿势。分红是直接打到王磊账上。回购是公司拿钱在市场上买回自己的股票然后注销，王磊手里的股数一股没变，但蛋糕被切成的片数少了，他那片自动变大。

\begin{egbox}{举个栗子}
三分甜净利 100 万，100 万股，股价 15 元，EPS 1 元。

它掏出 150 万现金，在市场上回购 10 万股，注销。股本变成 90 万股，利润还是 100 万，EPS 变成 100 万 ÷ 90 万 = \textbf{1.11 元}，涨了 11\%。

没多卖一杯奶茶，没多赚一分钱，EPS 涨了 11\%。

\end{egbox}

回购有两副面孔。股价低于内在价值时回购，等于替留下来的股东做了一笔划算买卖；股价高得离谱还借钱回购，那是拿股东的钱在高位接盘，顺便给 EPS 化个妆，好让管理层自己的期权顺利兑现。

税上回购确实占便宜。分红当期就要缴税，回购把收益留在股价里，不卖就不用交。A 股的股息税按持股期限差别化：持股超过 1 年暂免征收，1 个月以内全额计入应纳税所得额按 20\% 计征，中间那段减半。这条规则本身就在推着你拿久一点。

\section{称重机得有人真的去称}

格雷厄姆那句话被引用了八十年：市场短期是投票机，长期是称重机。

这句话有个从来没人明说的前提。\textbf{得有人真的在称。}

2025 到 2026 年，全球被动投资的规模第一次超过了主动管理。指数基金和 ETF 不研究公司，它们只干一件事：按指数权重买。某只股票被纳入指数，就必须买，不管贵不贵；被剔除，就必须卖，不管便宜不便宜。

王磊 2023 年也把大部分钱换成了沪深 300 指数基金，管理费从 1.5\% 降到 0.15\%，他觉得自己占了大便宜。他确实占了。

可如果所有人都这么想，谁来决定一家公司值多少钱？

\begin{deepbox}{进阶：可以跳过}
\textbf{一方说没问题。}定价是在边际上完成的。哪怕只剩 10\% 的资金还在做主动研究，只要它们能自由买卖，价格照样会被推回合理位置。指数基金按现有价格等比例买入，并不改变相对价格。何况被动化把管理费从 1.5\% 打到 0.05\%，这是实打实还给王磊们的钱。这一派背后是 Grossman 和 Stiglitz 那套老逻辑：只要主动管理还能赚到超额收益，就总有人来做，市场会自动维持一个刚好够用的研究量。

\textbf{另一方说有问题。}指数纳入和剔除本身就在制造跟基本面无关的价格波动，早被人抢跑套利。市值加权意味着涨得越多的股票权重越大，钱越往头部堆，流动性和估值一起挤在少数几家公司身上。下跌的时候被动资金的赎回同样机械，卖出不看估值，可能把踩踏放大。

\textbf{我的看法：}短期内价格发现不会失灵，主动资金的绝对规模还是太大了。但"钱按权重涌向头部"确实在改变市场结构，指数集中度处在历史高位是事实。至于被动化的比例到了多少会真正伤到价格发现，这件事目前没人知道答案，谁跟你说他知道，大概率是在卖产品。

\end{deepbox}

\section{王磊五年交了 4700 块过路费}

股市里的钱只有三个来源，性质完全不同。

第一种是企业盈利增长。三分甜今年比去年多赚了 40 万，所有股东按比例一起分，不需要任何人亏钱。这是股市长期回报唯一靠得住的来源。

第二种是估值变化。PE 从 15 涨到 20，公司一分钱没多赚，多出来那部分是后面接手的人付的。王磊赚的，正好是他多付的。

第三种是交易成本和税费。佣金、印花税、买卖价差、管理费。这部分钱谁也没赚走，它直接离开了池子。

王磊从 2021 年到 2026 年，账户从 1 万 2 加到 9 万 4，五年里买进卖出六十多次。佣金加印花税，一共交了 4700 多。

这 4700 块和大盘涨跌无关。涨也交，跌也交。

所以一群人在存量里高频倒手，整体一定是亏的。不是因为他们笨，是因为第三项永远带负号。第 20 章会把这个等式完整拆开。

\section{同一套打法，两个市场，两个结果}

\begin{tablewrap}
\begin{xltabular}{\linewidth}{>{\raggedright\arraybackslash}p{0.20\linewidth}XX}
\toprule
\bthead{维度} & \bthead{A 股} & \bthead{美股} \\
\midrule
\textbf{投资者构成} & 个人投资者贡献的成交额占比长期偏高 & 机构与被动资金主导，个人直接持股比例低得多 \\
\textbf{交易制度} & T+1；主板涨跌停 ±10\%，创业板与科创板 ±20\% & T+0，无涨跌幅限制，设有市场级熔断 \\
\textbf{退市} & 退市新规后节奏明显加快，历史存量仍在消化 & 常态化，每年大量公司因不达标或私有化退出 \\
\textbf{股息率} & 近年在政策推动下持续抬升，整体仍低于成熟市场 & 分红加回购的综合现金回报率更高 \\
\textbf{指数编制} & 上证综指覆盖全部上市股票，长期采用总股本加权 & 标普 500 由委员会筛选，自由流通市值加权，有盈利门槛 \\
\textbf{监管取向} & 融资功能与投资者保护并重，政策工具使用频繁 & 侧重信息披露与事后追责，集体诉讼是主要威慑 \\
\bottomrule
\end{xltabular}
\end{tablewrap}

王磊有个大学同学在纽约做量化，两人聊过好几次，发现同一套因子搬过去就失灵。

散户成交占比高的市场，情绪波动更大，短期反转更明显；退市不畅的市场，壳本身有价格，垃圾股会被反复炒；T+1 加涨跌停，会把今天没释放完的抛压原封不动推到明天开盘。

\begin{databox}{2026 数据}
2026 年 8 月 31 日收盘：上证指数 3986.30 点，当月涨 4.02\%；深证成指 14015 点；创业板指 3438.68 点；两市成交约 2.13 万亿元。中国公募基金规模约 38 万亿元（2026 年初）。日经 225 在 2026 年 6 月首次站上 70000 点。

多家机构判断，市场的核心逻辑正在从"估值修复"转向"盈利驱动"。这两个词对应的正好是前面那笔账：估值修复赚的是零和的那一块，盈利驱动赚的是正和的那一块。

\end{databox}

\begin{mythbox}{常见误解}
\mvfrow{误解}{"股价跌了一半，这公司肯定亏惨了。"}
\mvfrow{事实}{股价是股东之间转手的价钱，它根本不出现在公司的利润表里。三分甜股价从 18 跌到 9，它照样卖奶茶，照样赚它的 100 万。真正受影响的是间接那几条：再融资更贵、店长手里的期权变废纸、大股东质押可能爆仓、换股并购不划算。\textbf{公司亏钱是经营的事，股价下跌是定价的事。两者会互相拉扯，但不是一回事。}}
\mvfrow{误解}{"10 块钱的股票比 100 块的便宜。"}
\mvfrow{事实}{这句话是错的，错得很干净。股价是"每一片多少钱"，估值是"整个蛋糕多少钱"。A 公司切成 1 亿片，每片 10 元，市值 10 亿；B 公司切成 100 万片，每片 100 元，市值 1 亿。贵的那个反而便宜十倍。只看股价不看市值和盈利，等于问"这块蛋糕切一刀多少钱"。}
\mvfrow{误解}{"我买股票就是在支持实体经济。"}
\mvfrow{事实}{王磊那 1800 块给的是上一个股东，三分甜账上一分没多。真正把钱交到公司手里的，是参与 IPO 认购、参与定向增发、买公司新发行的债。\hlmark{但二级市场的存在，让公司能以更低成本增发、让创业者和早期投资人有地方退出、让期权有价值。}他不是在直接输血，他是在维护那套让输血成为可能的管道。这个作用是真的，只是隔了一层。}
\end{mythbox}

这本书不推荐任何标的，也不判断任何价格。上面所有的算术都是为了让你看懂这台机器怎么转。

\begin{takeaway}{本章一句话}
股票是剩余索取权：所有人拿完之后剩下的那份归你，所以它风险最大，也因此长期回报最高。

\begin{itemize}
  \item PE、PB 是贴现公式的简写，不是独立真理。周期股 PE 最低时往往最危险。
  \item 可持续增长率 ≈ ROE ×（1 − 分红率）。三分甜留下 68\% 的利润，它靠自己能长 13.6\%。
  \item 股市的钱只有三笔：盈利增长（正和）、估值变化（零和）、成本税费（负和）。王磊五年在第三笔上花掉 4700 块。
  \item 被动投资规模在 2025 到 2026 年首次超过主动管理。它会不会削弱价格发现，目前没人知道答案。
\end{itemize}

\end{takeaway}

\begin{bridgebox}
\textbf{下一章}股票和债券都用本国货币标价。可王磊要是想买德国的机器、日本的旅游、美国的芯片，他手里这串数字得先换成别人账本上的数字才算数。
\end{bridgebox}

